% =============================================================================
% Section 07 — Invariance and Composition
% Book I — Ontology of Viability
% Chapter 02 — Invariance
%
% Purpose:
%   To examine the relation between invariance and composition,
%   distinguishing the preservation of compositional determinations
%   from structural invariance and ontological identity.
%
% Conceptual Policy:
%   1. Compositional invariance is not identical to structural
%      invariance.
%   2. The persistence of components does not necessarily establish
%      the persistence of their composition.
%   3. Compositional change does not by itself establish the loss
%      of ontological identity.
%   4. The preservation of composition does not establish numerical
%      identity.
% =============================================================================

\section{Invariance and Composition}
\label{sec:invariance-and-composition}

\subsection{The Ontological Problem of Composition}

Composition concerns the manner in which multiple components
constitute or participate in the constitution of an existent.
It addresses not merely the presence of components, but also
their compositional arrangement and the determinate relations
through which they belong to a composed whole.

The relation between composition and invariance becomes
significant when an existent undergoes transformation while
some aspect of its composition remains preserved.

A composed existent may retain the same components while
their arrangement changes. It may also undergo the replacement
of some components while preserving a specified compositional
organization. In other cases, both components and their
compositional relations may change.

These possibilities indicate that compositional preservation
cannot be identified simply with the persistence of components
or with the persistence of the whole.

Compositional invariance must therefore be examined as a
specific form of invariance concerning the constitution of
an existent.

It does not, by itself, provide an exhaustive account of
the identity of that existent.

\subsection{Composition and the Presence of Components}

The presence of components is one aspect of composition,
but it is not sufficient to characterize composition as such.

A collection of components does not necessarily constitute
the same composed arrangement under every possible relation
or organization.

The same components may participate in different compositions.
Conversely, different components may occupy comparable
compositional positions within an organization.

This distinction requires separating two questions:

\begin{enumerate}
    \item Which components are present in a composed existent?
    \item How are those components compositionally constituted
          as part of that existent?
\end{enumerate}

The first question concerns component membership.
The second concerns compositional organization.

A claim of compositional invariance may concern either
of these respects or both. Consequently, the precise respect
of preservation must be specified.

The persistence of component membership cannot automatically
be treated as the persistence of composition. Likewise,
the replacement of a component does not necessarily imply
that every compositional determination has been altered.

\subsection{Compositional Invariance}

Compositional invariance refers to the preservation of a
specified compositional determination across a defined
transformation.

Such preservation may concern:

\begin{itemize}
    \item the membership of components within a composed whole;
    \item the number of components;
    \item the distribution of components among compositional
          positions;
    \item the relations through which components participate
          in a composition;
    \item the organization of components into a determinate
          compositional arrangement.
\end{itemize}

These respects are analytically distinguishable.
They need not be preserved together, and the preservation
of one does not establish the preservation of the others.

For example, the number of components may remain constant
while their membership changes. Component membership may
remain unchanged while their arrangement is reorganized.
A compositional arrangement may also remain equivalent
under a specified criterion even when particular components
are replaced.

Therefore, the assertion that composition remains invariant
requires a clear account of what aspect of composition
is being preserved.

Without such specification, compositional invariance
remains indeterminate.

\subsection{Component Replacement and Compositional Preservation}

An important case occurs when one or more components are
replaced while the broader compositional organization
remains preserved.

In this situation, the identity of individual components
does not necessarily determine the invariance of the
composition as a whole.

A replacement may preserve the compositional role or
position previously occupied by a component. It may also
preserve the number and organization of components while
altering their particular constitution.

However, the preservation of compositional organization
must not be assumed merely because replacement has occurred
without an immediately visible alteration of the whole.

Whether composition remains invariant depends on the
criterion by which the composition is being evaluated.

If component identity is constitutive of the specified
compositional determination, replacement may constitute
a loss of that invariance. If the relevant respect concerns
only a broader organization, replacement may be compatible
with its preservation.

Thus, component replacement does not have one universally
determined consequence for compositional invariance.

Its significance depends on the compositional determination
under consideration.

\subsection{Compositional Reorganization}

Composition may also change without the replacement
or disappearance of its components.

The same components may be reorganized into a different
compositional arrangement. Their distribution, participation,
or relations may change while their individual presence
remains preserved.

This possibility demonstrates that component persistence
and compositional invariance are not equivalent.

Compositional reorganization may occur at different levels.
A limited alteration may affect only one part of a composition,
while a more extensive transformation may alter the
organization of the composed whole.

The degree and significance of such transformation cannot
be determined solely by counting the components that
remain unchanged.

A compositional determination may be preserved despite
some local modifications, or it may cease to be preserved
even when all individual components remain present.

The assessment therefore requires an explicit criterion
of compositional preservation.

\subsection{Composition and Structural Invariance}

Compositional invariance and structural invariance are
closely related, but they are not identical.

Compositional invariance concerns the preservation of
a specified determination relating to the constitution
of a composed whole through its components.

Structural invariance concerns the preservation of a
specified organizational pattern or relational arrangement.

The two may coincide in particular cases. A compositional
arrangement may be preserved through the preservation
of its structural organization. Nevertheless, their
analytical scopes remain distinguishable.

A composition may preserve its component membership
while undergoing structural reorganization. Conversely,
a structural organization may remain equivalent even
when some of its components are replaced.

This distinction prevents compositional invariance
from being reduced to structural invariance.

It also prevents the assumption that every structural
transformation necessarily constitutes a transformation
of composition in every relevant respect.

The relationship between the two must be determined
according to the particular compositional and structural
determinations involved.

\subsection{Compositional Change and the Identity of the Existent}

The transformation of composition raises a further
ontological question: whether a change in the composition
of an existent entails a change in its identity.

No universal answer can be derived from compositional
invariance alone.

The replacement of a component, the reorganization
of components, or the alteration of compositional
relations may be relevant to the constitution of an
existent. However, none of these transformations,
considered in isolation, establishes that the existent
has ceased to be the existent that it is.

To infer the loss of identity from compositional change
would require an independent account of the ontological
significance of the particular composition under
consideration.

Likewise, the preservation of composition cannot
establish numerical identity. Distinct existents may
possess equivalent or indistinguishable compositional
determinations.

The same compositional arrangement may therefore
be instantiated by numerically distinct existents.

Compositional invariance is consequently insufficient
to establish either the persistence or the loss of
ontological identity.

\subsection{The Limits of Compositional Explanation}

Composition provides an important dimension for
describing the constitution of an existent. Nevertheless,
it cannot be treated as an exhaustive explanation
of that existent.

An existent is not necessarily reducible to the sum
of its components, nor can its identity be established
solely through the preservation of their membership
or arrangement.

Such reduction would assume that every ontologically
relevant determination of an existent is contained
within its composition.

That assumption cannot be established through the
analysis of compositional invariance itself.

The ontological significance of composition must
therefore be acknowledged without granting it
explanatory exclusivity.

Compositional determinations may contribute to the
constitution and characterization of an existent,
but their preservation or transformation does not
independently settle the question of its identity.

\subsection{Synthesis}

The analysis establishes that compositional invariance
concerns the preservation of a specified compositional
determination across a defined transformation.

Its assessment requires distinguishing:

\begin{itemize}
    \item component membership from compositional organization;
    \item component persistence from compositional preservation;
    \item component replacement from the transformation
          of the entire composition;
    \item compositional invariance from structural invariance;
    \item compositional equivalence from numerical identity;
    \item compositional transformation from the loss
          of ontological identity.
\end{itemize}

The preservation of composition does not necessarily
establish the persistence of an existent, and the
transformation of composition does not necessarily
establish its loss of identity.

Compositional invariance must therefore remain a
determinate and limited category within the broader
analysis of invariance.

It identifies a possible mode of preservation in
the constitution of an existent, without reducing
the identity of that existent to the persistence
of its components or their compositional arrangement.

% =============================================================================
% End of Section 07
% =============================================================================